% !TEX program = xelatex

\documentclass[11pt,a4paper]{article}
\usepackage{geometry}
\geometry{margin=2cm}
\usepackage{array}
\usepackage{booktabs}
\usepackage{fontspec}
\usepackage{arabxetex}
\usepackage{xcolor}
\usepackage{titlesec}

% Your System Font Pattern
\newfontfamily\arabicfont[Script=Arabic]{Amiri}
\newcommand{\ar}[1]{{\arabicfont #1}}

% Design Formatting Elements
\definecolor{primary}{HTML}{1A365D}
\definecolor{secondary}{HTML}{2B6CB0}
\definecolor{accent}{HTML}{C53030}

\titleformat{\section}{\color{primary}\normalfont\large\bfseries}{\thesection}{1em}{}[{\titrule[0.8pt]}]
\titleformat{\subsection}{\color{secondary}\normalfont\normalsize\bfseries}{\thesubsection}{1em}{}

\setlength{\parindent}{0pt}
\setlength{\parskip}{0.8em}

\begin{document}

\begin{center}
    {\color{primary}\LARGE\bfseries Comprehensive Reference: Arabic Case Endings} \\
    \vspace{0.3cm}
    {\color{secondary}\large A Structural Blueprint of Singular, Dual, and Broken Plural Forms} \\
    \vspace{0.2cm}
    \strut\hrule width \linewidth height 1.5pt \relax
\end{center}

\vspace{0.5cm}

\section{The Core Principle of Arabic I'rāb (Vocalization)}
In Classical and Quranic Arabic, the final vowel or letter configuration of a noun changes systematically depending on its structural role within a sentence:

\begin{itemize}
    \item \textbf{\color{primary}Nominative Case [ \ar{مَرْفُوع} -- Marfūʿ ]:} Typically marked by a \textbf{\color{primary}ḍamma [ \ar{ـُ} ]} on singular nouns. It indicates the Subject or Doer performing the action.
    \item \textbf{\color{secondary}Accusative Case [ \ar{مَنْصُوب} -- Manṣūb ]:} Typically marked by a \textbf{\color{secondary}fatḥa [ \ar{ـَ} ]} on singular nouns. It identifies the Direct Object receiving the action.
    \item \textbf{\color{accent}Genitive Case [ \ar{مَجْرُور} -- Majrūr ]:} Typically marked by a \textbf{\color{accent}kasra [ \ar{ـِ} ]} on singular nouns. It specifies possession or follows a grammatical preposition.
\end{itemize}

\section{Illustrative Sentences: The Three Primary Cases}

\subsection{2.1 Nominative $\rightarrow$ The Doer}
\textbf{Sentence Structure:} Verb + Doer + Object \\
\textbf{Example Text:} \hfill {\Large \textarabic{خَلَقَ اللَّهُ الإِنْسَانَ}} \\
\textbf{Transliteration:} \textit{Khalaqa Allāhu al-insāna} \\
\textbf{Syntactic Breakdown:}
\begin{itemize}
    \item \textbf{Verb:} \textarabic{خَلَقَ} (\textit{Khalaqa}) $\rightarrow$ Created
    \item \textbf{Doer / Subject:} \textarabic{اللَّهُ} (\textit{Allāhu}) $\rightarrow$ Marked by the final ḍamma
    \item \textbf{Direct Object:} \textarabic{الإِنْسَانَ} (\textit{al-insāna}) $\rightarrow$ Mankind
\end{itemize}
\textbf{Translation:} ``Allah created mankind.''

\subsection{2.2 Accusative $\rightarrow$ The Receiver}
\textbf{Example Text:} \hfill {\Large \textarabic{نَصَرَ الرَّجُلُ يُوسُفَ}} \\
\textbf{Transliteration:} \textit{Naṣara ar-rajulu Yūsufa} \\
\textbf{Syntactic Breakdown:}
\begin{itemize}
    \item \textbf{Verb:} \textarabic{نَصَرَ} (\textit{Naṣara}) $\rightarrow$ Helped
    \item \textbf{Doer / Subject:} \textarabic{الرَّجُلُ} (\textit{ar-rajulu}) $\rightarrow$ The man (Marked by ḍamma)
    \item \textbf{Direct Object / Receiver:} \textarabic{يُوسُفَ} (\textit{Yūsufa}) $\rightarrow$ Joseph (Marked by the final fatḥa)
\end{itemize}
\textbf{Translation:} ``The man helped Joseph.''

\subsection{2.3 Genitive $\rightarrow$ Possession or Prepositional Object}
\textbf{Example A (Possession):} \hfill {\Large \textarabic{كِتَابُ الرَّجُلِ}} \\
\textbf{Transliteration:} \textit{Kitābu ar-rajuli} \\
\textbf{Meaning:} ``The man's book'' / ``The book of the man.'' (Marked by the final kasra)

\textbf{Example B (After a Preposition):} \hfill {\Large \textarabic{ذَهَبْتُ إِلَى الْمَسْجِدِ}} \\
\textbf{Transliteration:} \textit{Dhahabtu ilā al-masjidi} \\
\textbf{Syntactic Breakdown:}
\begin{itemize}
    \item \textbf{Phrase:} \textarabic{ذَهَبْتُ} (\textit{Dhahabtu}) $\rightarrow$ I went
    \item \textbf{Preposition:} \textarabic{إِلَى} (\textit{ilā}) $\rightarrow$ To
    \item \textbf{Prepositional Object:} \textarabic{الْمَسْجِدِ} (\textit{al-masjidi}) $\rightarrow$ The mosque (Marked by the final kasra)
\end{itemize}
\textbf{Translation:} ``I went to the mosque.''

\pagebreak

\section{The Complete Declension Framework for the Root: \ar{ر - ج - ل}}
\begin{table}[h!]
\centering
\small
\renewcommand{\arraystretch{1.4}}
\begin{tabular}{lllll}
\toprule
\textbf{\color{primary}Grammatical Number} & \textbf{\color{primary}Case} & \textbf{\color{primary}Arabic Form} & \textbf{\color{primary}Transliteration} & \textbf{\color{primary}Primary Syntactic Function} \\ 
\midrule
\textbf{Singular} & Nominative & {\Large \textarabic{الرَّجُلُ}} & \textit{ar-rajulu} & Sentence Subject / Doer \\
(One Man) & Accusative & {\Large \textarabic{الرَّجُلَ}} & \textit{ar-rajula} & Direct Object / Receiver \\
 & Genitive & {\Large \textarabic{الرَّجُلِ}} & \textit{ar-rajuli} & Post-Preposition / Possession \\ 
\midrule
\textbf{Dual} & Nominative & {\Large \textarabic{الرَّجُلَانِ}} & \textit{ar-rajulāni} & Sentence Subject / Doer \\
(Two Men) & Accusative & {\Large \textarabic{الرَّجُلَيْنِ}} & \textit{ar-rajulayni} & Direct Object / Receiver \\
 & Genitive & {\Large \textarabic{الرَّجُلَيْنِ}} & \textit{ar-rajulayni} & Post-Preposition / Possession \\ 
\midrule
\textbf{Broken Plural} & Nominative & {\Large \textarabic{الرِّجَالُ}} & \textit{ar-rijālu} & Sentence Subject / Doer \\
(Three or More Men) & Accusative & {\Large \textarabic{الرِّجَالَ}} & \textit{ar-rijāla} & Direct Object / Receiver \\
 & Genitive & {\Large \textarabic{الرِّجَالِ}} & \textit{ar-rijāli} & Post-Preposition / Possession \\ 
\bottomrule
\end{tabular}
\caption{Systematic Declinational Blueprint of the Noun: \ar{رَجُل}}
\end{table}

\section{Structural Contrast: The Vowel Swap Experiment}
Shifting the short vowels re-allocates semantic meaning entirely without modifying word order:

\vspace{0.3cm}

\begin{center}
\begin{minipage}{0.48\textwidth}
    \centering
    \textbf{Historical Quranic Fact} \\[0.5em]
    {\Large \textarabic{قَتَلَ دَاوُدُ جَالُوتَ}} \\[0.5em]
    \small \textit{Qatala Dāwūdu Jālūta}
    \begin{itemize}
        \item[] \textarabic{دَاوُدُ} $\rightarrow$ The Doer / Subject (David / marked by ḍamma)
        \item[] \textarabic{جَالُوتَ} $\rightarrow$ The Object (Goliath / marked by fatḥa)
    \end{itemize}
    \vspace{0.5em}
    \textbf{Meaning:} ``David killed Goliath.''
\end{minipage}
\hfill
\begin{minipage}{0.48\textwidth}
    \centering
    \textbf{Inverted Semantic Paradox} \\[0.5em]
    {\Large \textarabic{قَتَلَ دَاوُدَ جَالُوتُ}} \\[0.5em]
    \small \textit{Qatala Dāwūda Jālūtu}
    \begin{itemize}
        \item[] \textarabic{دَاوُدَ} $\rightarrow$ The Object (David / marked by fatḥa)
        \item[] \textarabic{جَالُوتُ} $\rightarrow$ The Doer / Subject (Goliath / marked by ḍamma)
    \end{itemize}
    \vspace{0.5em}
    \textbf{Meaning:} ``Goliath killed David.''
\end{minipage}
\end{center}

\vspace{1.5cm}
\begin{center}
    \small\color{gray} This document serves as an analytical reference guide for Classical Arabic morphology and syntax grammar tracking.
\end{center}

\end{document}
